\chapter{Quadro teorico}
\label{chap:quadro-teorico}
\label{cap:quadro_teorico}

\section{People Analytics e turnover}

\subsection{Definizione ed evoluzione della People Analytics}

La \textit{People Analytics}, nota anche come \textit{Talent Analytics} o \textit{Workforce Analytics}, può essere definita come l'integrazione sistematica di fonti di dati eterogenee, provenienti sia dall'interno sia dall'esterno dell'organizzazione, finalizzata a rispondere a domande strategiche di business relative al capitale umano di un'impresa. A differenza della reportistica HR tradizionale, che si limita a descrivere ciò che è già accaduto, la People Analytics si concentra sulla previsione dei risultati futuri e sulla creazione di valore strategico attraverso l'analisi sistematica del comportamento dei dipendenti \cite{Isson_Harriott_2016}.

L'interesse verso la misurazione sistematica del lavoro umano non costituisce un fenomeno recente. Le radici storiche della disciplina possono essere rintracciate all'inizio del XX secolo, con lo \textit{Scientific Management} di Frederick Taylor (1911) e le prime applicazioni della psicologia industriale a opera di Hugo Münsterberg (1913), a cui si affiancò, durante la Prima Guerra Mondiale, l'adozione dei primi test mentali su larga scala per la selezione del personale militare. Un passaggio decisivo verso la quantificazione delle risorse umane si è avuto alla fine degli anni Settanta, quando Jac Fitz-enz propose le prime metriche formali per l'HR, ponendo le basi per lo sviluppo del benchmarking di settore. Negli anni Novanta, l'adozione della \textit{Balanced Scorecard} e della sua successiva declinazione HR rese esplicito il legame tra pratiche HR e performance organizzativa, mentre a partire dalla metà degli anni Duemila la diffusione dei \textit{Talent Management Systems} (TMS) ha consentito di automatizzare i processi HR e di archiviare in modo strutturato i dati sui dipendenti. Soltanto intorno al 2010, infine, l'analisi predittiva ha iniziato a consolidarsi come pratica effettiva nei dipartimenti delle risorse umane, nell'ambito di un più ampio spostamento verso un'integrazione multidisciplinare di fonti dati allineate agli obiettivi strategici del vertice aziendale.

L'evoluzione contemporanea della disciplina è sospinta da un insieme di forze convergenti. Da un lato, l'esplosione dei \textit{big data} e la diffusione di sorgenti digitali relative al lavoro hanno reso disponibile una quantità di informazioni sul talento senza precedenti, trasformandole in una potenziale leva competitiva. Dall'altro, la progressiva riduzione della permanenza media dei dipendenti presso la medesima organizzazione ha accresciuto l'urgenza di strumenti capaci di cogliere precocemente i segnali di un possibile abbandono. A queste dinamiche si aggiunge il mutato atteggiamento delle generazioni più giovani nei confronti del lavoro, che richiede forme più continue e strutturate di monitoraggio dell'\textit{engagement}.

In questo scenario, la funzione HR si trova coinvolta in un progressivo spostamento da un paradigma tradizionale, prevalentemente reattivo e orientato alla dimensione amministrativa e \textit{soft} della gestione del personale, a un paradigma emergente in cui la gestione del talento identifica in modo proattivo rischi e opportunità legati al capitale umano, adottando un approccio basato sull'evidenza empirica. È in questa prospettiva che la People Analytics è stata definita il ``\textit{Moneyball} delle risorse umane'', in analogia con la trasformazione \textit{data-driven} che ha investito ambiti professionali inizialmente poco permeabili all'uso strategico dei dati.

\subsection{Fattori di attrition}

L'\textit{attrition}, ovvero l'uscita volontaria dei dipendenti dall'organizzazione, può essere ricondotta a un'interazione complessa tra condizioni di lavoro interne, pressioni del mercato esterno e caratteristiche individuali del lavoratore. La rassegna della letteratura consente di raggruppare i principali antecedenti in quattro macro-aree, utili a delineare lo spazio delle variabili su cui un modello predittivo di People Analytics è chiamato a operare \cite{Isson_Harriott_2016}.

Un primo gruppo di determinanti attiene all'ambiente di lavoro e alle pratiche gestionali. Un carico di lavoro eccessivo, unito alla progressiva erosione dei periodi di recupero, tende ad accrescere lo stress percepito e ad aumentare sensibilmente la probabilità di abbandono nei mesi successivi. Altrettanto rilevante è la qualità della relazione con il \textit{management}: il comportamento dei superiori diretti, in particolare all'interno di progetti di grandi dimensioni, incide in modo profondo sull'esperienza lavorativa quotidiana e, di conseguenza, sulla permanenza organizzativa. Anche la scarsità di \textit{feedback} strutturato sulle prestazioni e la mancanza di un riconoscimento tangibile dei risultati alimentano un progressivo disallineamento tra aspettative e percezioni, contribuendo a erodere il legame con l'organizzazione.

Un secondo gruppo di fattori riguarda la dimensione economica e lo sviluppo professionale. La percezione di un'inadeguata equità retributiva e l'assenza di prospettive credibili di crescita interna figurano fra le ragioni più frequentemente richiamate dai lavoratori che scelgono di cambiare datore di lavoro. A queste si affiancano determinanti di ordine culturale e generazionale: la letteratura sottolinea come le coorti più giovani attribuiscano un valore crescente al significato personale del lavoro, all'allineamento con la \textit{mission} aziendale, al senso di appartenenza e alla responsabilità sociale dell'impresa \cite{deloitte2026trends}. In assenza di una missione riconoscibile e condivisa, tali lavoratori mostrano una maggiore propensione alla mobilità. In questa stessa chiave vanno letti i mutamenti più recenti legati alla diffusione dell'intelligenza artificiale nei processi di lavoro, che ha introdotto nuovi fattori di rischio quali l'aumento dei carichi cognitivi, la riduzione dell'autonomia, un senso di isolamento dovuto alla minore collaborazione fra colleghi e l'incertezza sulla stabilità futura del proprio ruolo \cite{deloitte2026trends}.

Una terza famiglia di antecedenti ha natura strettamente individuale. Tratti di personalità e \textit{fit} culturale con l'organizzazione risultano fra i predittori più significativi della permanenza: una marcata discrepanza fra i valori personali del dipendente e quelli dell'impresa rappresenta infatti una delle cause più ricorrenti di \textit{turnover} precoce. Simili variabili individuali sono notoriamente difficili da osservare per via diretta e, quando non disponibili, il loro effetto tende a essere parzialmente assorbito dalle variabili comportamentali e contrattuali registrate nei sistemi informativi HR.

Un ultimo gruppo di fattori, infine, opera dall'esterno dell'organizzazione. La presenza di mercati del lavoro competitivi, con bassi tassi di disoccupazione per i profili qualificati e offerte retributive aggressive da parte dei concorrenti, agisce come potente leva di attrazione verso l'esterno. La progressiva diffusione delle piattaforme professionali digitali, rendendo immediatamente accessibili nuove opportunità e notifiche personalizzate, ha inoltre reso la forza lavoro più reattiva ai segnali di mercato, accorciando i tempi di decisione che precedono le dimissioni.

Un elemento importante ai fini del presente lavoro riguarda l'orizzonte temporale entro cui tali segnali iniziano a manifestarsi. Le evidenze disponibili suggeriscono che i primi indicatori di disimpegno compaiano in media con un anticipo compreso fra i nove e i dodici mesi rispetto alle dimissioni effettive. È proprio in questa finestra temporale che l'impiego di modelli predittivi di People Analytics assume rilevanza strategica, consentendo di identificare i dipendenti a rischio con sufficiente anticipo per attivare interventi di \textit{retention} mirati.

\subsection{Limiti degli approcci tradizionali e vantaggi dei dati strutturati}

Nonostante la crescente disponibilità di dati sul capitale umano, una parte significativa delle decisioni relative alla gestione dei talenti continua a fondarsi sull'intuizione personale, sull'esperienza pregressa e su relazioni interpersonali difficilmente oggettivabili. La letteratura evidenzia come tale approccio, definito talvolta ``basato sull'esperienza educata'', risulti privo della profondità analitica necessaria per orientare scelte strategiche in contesti organizzativi complessi \cite{Isson_Harriott_2016}.

A questa fragilità strutturale si accompagnano ulteriori limiti. I metodi tradizionali di selezione, come le interviste non strutturate, hanno mostrato un potere predittivo sulla performance lavorativa futura prossimo a quello di una scelta casuale; analogamente, criteri storicamente consolidati quali il voto di laurea o i punteggi dei test attitudinali si sono rivelati indicatori poco affidabili del rendimento effettivo. Sul versante della \textit{retention}, gli strumenti convenzionali di ascolto organizzativo --- quali le \textit{survey} periodiche e le \textit{exit interview} --- operano in modo prevalentemente retrospettivo, registrando il disagio soltanto quando le dimissioni si sono già verificate o sono imminenti, senza offrire margini per un intervento tempestivo. Anche i canali di reclutamento più diffusi, dalle bacheche di annunci alle fiere universitarie, sono divenuti risorse generiche, accessibili a tutti i concorrenti e pertanto incapaci di conferire un vantaggio competitivo distintivo. In un simile contesto, la funzione HR rischia di essere percepita come un centro di costo a vocazione prevalentemente amministrativa e normativa, distante dalle priorità strategiche del vertice aziendale.

L'adozione di dati strutturati e di tecniche analitiche avanzate consente di superare questi limiti e di riposizionare la gestione delle risorse umane come funzione strategica. Il beneficio principale risiede nella possibilità di transitare da decisioni basate sull'aneddoto a strategie fondate su evidenze analitiche, legittimando l'HR come interlocutore credibile del \textit{top management}. In secondo luogo, l'integrazione di fonti dati eterogenee --- dalle rilevazioni interne ai dati di mercato del lavoro, dall'analisi dei flussi digitali alle metriche di \textit{engagement} --- permette di rispondere a domande aziendali complesse in modo rapido e accurato, costruendo modelli predittivi in grado di identificare in anticipo i dipendenti a rischio di abbandono e di stimare il ritorno sull'investimento del capitale umano. A titolo illustrativo, la letteratura riporta il caso di un'organizzazione che, sostituendo i criteri di selezione basati sull'esperienza con variabili legate ai tratti di personalità, ha ottenuto una riduzione significativa del tasso di \textit{attrition}. Infine, la disponibilità di dati longitudinali rende possibile l'applicazione al talento di modelli ispirati alla gestione del ciclo di vita del cliente, tracciando ciascun dipendente dal primo contatto con l'organizzazione fino all'eventuale uscita, così da ottimizzare le strategie di \textit{retention} lungo l'intero percorso professionale.

È in questa transizione dall'intuizione all'evidenza che si colloca il contributo della People Analytics: una ``scienza delle decisioni'' applicata alle persone, analoga a quanto già avvenuto nella finanza e nel marketing.

\section{Explainable AI (XAI)}

\subsection{Definizione di interpretabilità}

L'interpretabilità nel \textit{machine learning} è un concetto che non possiede una definizione matematica univoca, ma viene descritto attraverso prospettive complementari nella letteratura \cite{molnar2025interpretable}. Una prima accezione la identifica con il grado in cui un essere umano è in grado di comprendere la causa di una decisione prodotta da un modello. Una seconda prospettiva, di natura più operativa, considera un metodo interpretabile quando un utente riesce a prevedere correttamente ed efficientemente i risultati del metodo stesso: in altri termini, quanto più il comportamento del modello risulta anticipabile, tanto maggiore è la sua interpretabilità. Una terza definizione descrive l'interpretabilità come l'atto di mappare concetti astratti, generati internamente dal modello, in una forma che sia comprensibile per l'uomo \cite{roscher2020explainable}. Infine, in un'accezione più ampia, il termine può essere inteso come un concetto ombrello che racchiude l'estrazione di conoscenze rilevanti da un modello di \textit{machine learning}, sia per quanto riguarda le relazioni contenute nei dati, sia per quelle apprese dal modello durante l'addestramento \cite{molnar2025interpretable}.

In termini comparativi, un modello si considera più interpretabile di un altro quando le sue decisioni risultano più facili da comprendere per un essere umano. Questa definizione relativa riflette il fatto che l'interpretabilità non è una proprietà binaria, bensì un continuum lungo il quale i modelli si collocano in posizioni differenti. L'obiettivo ultimo dell'interpretabilità è consentire a chi interagisce con il modello di comprendere il \textit{perché} di determinate decisioni o previsioni, facilitando così la risoluzione di problemi che derivano da una formalizzazione necessariamente incompleta dei compiti del mondo reale \cite{molnar2025interpretable}.

\subsection{Interpretabilità ed esplicabilità}

I termini \textit{interpretability} ed \textit{explainability} vengono frequentemente impiegati come sinonimi nella letteratura sul \textit{machine learning}, tuttavia tra essi sussistono sfumature tecniche che meritano di essere esplicitate, nonostante l'assenza di un accordo univoco nella comunità scientifica \cite{molnar2025interpretable}.

Secondo la distinzione proposta da Roscher et al., l'\textit{interpretability} riguarda la mappatura di un concetto astratto prodotto dal modello in una forma comprensibile per l'essere umano, mentre l'\textit{explainability} rappresenta un requisito più forte, che presuppone l'interpretabilità ma richiede in aggiunta l'inserimento di un contesto supplementare capace di rendere la spiegazione pienamente significativa per il destinatario \cite{roscher2020explainable}. In questa prospettiva, l'esplicabilità compie un passo ulteriore rispetto all'interpretabilità, collocando la comprensione del modello all'interno di un quadro di riferimento più ampio.

Sul piano terminologico, il termine \textit{explanation} (spiegazione) viene generalmente riservato ai metodi locali, il cui scopo è rendere conto di una singola previsione specifica, piuttosto che del comportamento complessivo del modello \cite{molnar2025interpretable}. Nella pratica, tuttavia, la distinzione tra i due concetti tende a sfumare. L'approccio pragmatico suggerito dalla letteratura consiste nel trattarli come un termine ombrello che racchiude l'insieme delle tecniche finalizzate all'estrazione di conoscenza rilevante da un modello, sia riguardo alle relazioni presenti nei dati sia a quelle apprese durante l'addestramento \cite{molnar2025interpretable}.

\subsection{Caratteristiche di una buona spiegazione}

La ricerca nelle scienze cognitive e sociali ha messo in luce che una spiegazione efficace per l'essere umano non coincide necessariamente con una descrizione completa e tecnicamente esaustiva di tutte le cause di un fenomeno, ma deve possedere proprietà specifiche che riflettono il modo in cui le persone elaborano e valutano le informazioni \cite{molnar2025interpretable}.

La prima proprietà è la \textbf{contrastività}. Gli esseri umani raramente si interrogano sulle cause di un evento in termini assoluti; tendono piuttosto a chiedersi perché sia accaduto un certo esito \textit{invece di un altro}. Il ragionamento è intrinsecamente controfattuale: chi riceve il rifiuto di una richiesta non desidera conoscere tutti i fattori generali che hanno condotto alla decisione, bensì quelli che avrebbero dovuto essere diversi affinché l'esito fosse favorevole.

La seconda proprietà è la \textbf{selettività}. Le persone non si aspettano un elenco completo di tutte le cause che hanno contribuito a un risultato, ma preferiscono spiegazioni brevi che individuino al massimo due o tre fattori principali, anche quando la realtà sottostante è considerevolmente più complessa.

La terza proprietà riguarda la dimensione \textbf{sociale} della spiegazione. Ogni spiegazione si inserisce in un'interazione tra chi la formula e chi la riceve: il suo contenuto, il livello di dettaglio e il registro comunicativo devono adattarsi al destinatario e al contesto in cui avviene lo scambio. Una spiegazione rivolta a un esperto del dominio differirà in modo sostanziale da quella destinata a un interlocutore privo di competenze tecniche.

La quarta proprietà è la \textbf{focalizzazione sull'anormale}. Gli esseri umani tendono a considerare particolarmente convincenti le spiegazioni che si concentrano su cause rare o insolite. Quando un evento inatteso ha influenzato il risultato, esso viene percepito come la causa principale, anche qualora altri fattori ordinari abbiano esercitato un peso analogo.

La quinta proprietà è la \textbf{veridicità}: una spiegazione dovrebbe essere il più possibile fedele alla realtà. Tuttavia, nella valutazione umana, la selettività e il contrasto risultano spesso più importanti della completezza. Un'approssimazione parziale ma chiara viene generalmente preferita a una ricostruzione esaustiva ma incomprensibile.

La sesta proprietà è la \textbf{coerenza con le credenze precedenti}. Le persone tendono a svalutare o ignorare le spiegazioni che contraddicono le proprie convinzioni pregresse, un fenomeno riconducibile al \textit{bias} di conferma. Una spiegazione efficace è pertanto quella che si armonizza con il modello mentale e le conoscenze pregresse del destinatario, pur introducendo elementi informativi nuovi.

Infine, in assenza di eventi anormali, vengono considerate buone le spiegazioni \textbf{generali}, ovvero quelle applicabili a molti casi simili e fondate su relazioni probabilisticamente solide.

Queste proprietà hanno implicazioni dirette per la progettazione di metodi di \textit{Explainable AI}: affinché una spiegazione algoritmica risulti realmente utile, non è sufficiente che sia accurata dal punto di vista del modello, ma deve essere formulata in modo da rispettare i meccanismi cognitivi e sociali attraverso i quali gli esseri umani comprendono e valutano le informazioni \cite{molnar2025interpretable}.

\subsection{Quando l'interpretabilità è necessaria e quando no}

L'esigenza di interpretabilità nel \textit{machine learning} nasce dal fatto che una singola metrica di performance, come l'accuratezza, fornisce una descrizione incompleta della maggior parte dei compiti del mondo reale. Per affrontare problemi complessi, non è sufficiente conoscere \textit{cosa} è stato previsto: è spesso indispensabile comprendere il \textit{perché} \cite{molnar2025interpretable}.

Le ragioni per cui l'interpretabilità risulta necessaria possono essere ricondotte a sei ambiti principali. Il primo riguarda il \textbf{miglioramento e il \textit{debugging} del modello}: l'interpretabilità consente di identificare se un modello ha appreso scorciatoie spurie --- i cosiddetti \textit{Clever Hans predictors} --- ovvero se si affida a caratteristiche non causali per formulare le proprie previsioni. Essa permette inoltre di individuare errori nell'ingegneria delle \textit{feature} e di generare ipotesi per il loro perfezionamento.

Il secondo ambito concerne la \textbf{giustificazione e la responsabilità}. In numerosi settori, fra cui quello bancario e quello sanitario, è necessario fornire agli \textit{stakeholder} una giustificazione delle decisioni algoritmiche. L'interpretabilità consente di contestare decisioni errate e di verificare che il modello operi in conformità con le normative vigenti; in ambito europeo, il Regolamento generale sulla protezione dei dati (GDPR) riconosce il diritto dell'interessato a non essere sottoposto a decisioni basate unicamente su un trattamento automatizzato e a ottenere informazioni significative sulla logica utilizzata \cite{GDPR_Art22}.

Il terzo ambito attiene alla \textbf{sicurezza e alla gestione del rischio}. In contesti ad alto impatto è essenziale accertarsi che il sistema abbia appreso rappresentazioni prive di errori sistematici; comprendere il processo decisionale del modello aiuta a prevedere casi limite e a garantire la robustezza complessiva del sistema.

Il quarto ambito riguarda il \textbf{rilevamento dei \textit{bias} e la \textit{fairness}}. I modelli possono ereditare pregiudizi dai dati di addestramento, producendo decisioni discriminatorie nei confronti di gruppi sottorappresentati. L'interpretabilità funge da strumento di \textit{debugging} per individuare distorsioni che non verrebbero rilevate dalla sola funzione di perdita.

Il quinto ambito è quello della \textbf{scoperta scientifica}. Numerose discipline impiegano il \textit{machine learning} per estrarre conoscenze da grandi insiemi di dati; l'interpretabilità consente di trasformare il modello stesso in una fonte di comprensione dei fenomeni, rendendo esplicite le relazioni tra le variabili.

Il sesto ambito, infine, riguarda la \textbf{curiosità umana e l'accettazione sociale}. Gli esseri umani possiedono il desiderio di trovare un significato nel mondo e di armonizzare le contraddizioni; un sistema capace di spiegare il proprio comportamento favorisce la fiducia e facilita l'integrazione degli algoritmi nella pratica quotidiana \cite{molnar2025interpretable}.

Esistono tuttavia situazioni in cui l'interpretabilità non è necessaria o può risultare persino controproducente. In \textbf{ambienti a basso rischio}, dove un errore del modello non comporta conseguenze rilevanti, l'assenza di spiegazioni non costituisce un problema. Analogamente, per \textbf{problemi ampiamente studiati e consolidati}, nei quali si è accumulata sufficiente esperienza pratica e il sistema opera in modo affidabile, il beneficio marginale dell'interpretabilità è limitato. In alcuni casi, l'interpretabilità può essere intenzionalmente evitata per \textbf{prevenire la manipolazione del sistema}: se i criteri decisionali di un modello fossero interamente trasparenti, gli utenti potrebbero alterare strategicamente i propri dati per ottenere un risultato favorevole, senza che la loro condizione reale sia mutata. Questo rischio è particolarmente concreto quando il modello si basa su variabili \textit{proxy} che correlano con il risultato senza esserne la causa diretta \cite{molnar2025interpretable}.

In sintesi, la necessità di interpretabilità emerge quando la formalizzazione del problema è incompleta e la sola previsione non è sufficiente a risolvere pienamente il compito nel mondo reale. Se il compito è ben definito, sicuro e privo di implicazioni rilevanti per le persone coinvolte, l'interpretabilità può legittimamente passare in secondo piano.

\subsection{L'effetto Rashomon}

L'effetto Rashomon nel \textit{machine learning} designa la situazione in cui esistono molteplici modelli che presentano prestazioni predittive comparabili, ma forniscono interpretazioni diverse e potenzialmente incompatibili dei dati sottostanti \cite{molnar2025interpretable}. Il termine trae origine dall'omonimo film di Akira Kurosawa (1950), nel quale quattro testimoni narrano versioni differenti del medesimo evento: ciascuna ricostruzione spiega i fatti in modo altrettanto plausibile, eppure le versioni risultano inconciliabili fra loro.

In ambito modellistico, questo fenomeno solleva sfide rilevanti per l'interpretabilità. In primo luogo, la \textbf{molteplicità dei modelli} implica che diversi algoritmi o configurazioni possano raggiungere la medesima accuratezza basandosi tuttavia su relazioni differenti tra le \textit{feature} e la variabile obiettivo. In secondo luogo, tale molteplicità genera un'\textbf{incertezza nell'interpretazione}: se due modelli funzionano in modo equivalente dal punto di vista predittivo, ma raccontano storie diverse su come i dati interagiscono, non è chiaro quale di queste rifletta in modo più fedele la realtà del fenomeno osservato. Infine, l'effetto Rashomon complica la \textbf{scoperta di conoscenze}: per utilizzare un modello come surrogato della realtà, è necessario assumere che la sua struttura rispecchi il mondo; l'esistenza di modelli alternativi ed equivalenti mina questa assunzione, poiché ciascuno di essi potrebbe suggerire conclusioni diverse sullo stesso fenomeno \cite{molnar2025interpretable}.

L'effetto Rashomon evidenzia dunque come una buona performance predittiva non garantisca un'unica spiegazione corretta, complicando il compito di chi cerca di comprendere i meccanismi sottostanti ai dati attraverso l'analisi dei modelli. Questa considerazione è particolarmente rilevante nel contesto della presente tesi, dove tre diversi framework di \textit{Explainable AI} vengono applicati al medesimo modello: le eventuali discordanze tra le spiegazioni prodotte non devono essere interpretate necessariamente come un limite degli strumenti, ma possono riflettere la pluralità intrinseca delle interpretazioni compatibili con le stesse evidenze empiriche.

\subsection{SHAP (SHapley Additive exPlanations)}

Il framework SHAP (\textit{SHapley Additive exPlanations}), proposto da Lundberg e Lee, rappresenta un approccio unificato all'interpretazione delle previsioni dei modelli di \textit{machine learning}, fondato sui valori di Shapley della teoria dei giochi cooperativi \cite{lundberg2017unified}.

\subsubsection{Fondamenti nella teoria dei giochi cooperativi}

I valori di Shapley, originariamente formulati nell'ambito della teoria dei giochi cooperativi, risolvono il problema dell'attribuzione equa del merito tra i partecipanti a un gioco collaborativo. SHAP adatta questo concetto al contesto predittivo, rappresentando ogni previsione del modello come un gioco in cui le \textit{feature} sono i giocatori che cooperano per spostare la previsione dal valore medio di base (la media delle previsioni sull'intero dataset) al valore finale calcolato per una specifica istanza \cite{lundberg2017unified}.

Per ciascuna previsione, SHAP assegna a ogni \textit{feature} un valore di attribuzione $\phi_i$ che quantifica il contributo marginale di quella \textit{feature} al risultato. Il calcolo si basa sul cambiamento dell'aspettativa condizionata del modello quando la \textit{feature} viene inclusa in un dato sottoinsieme: concretamente, il valore di Shapley di una \textit{feature} corrisponde alla media ponderata dei suoi contributi marginali, calcolata su tutte le possibili combinazioni di \textit{feature} \cite{lundberg2017unified}.

\subsubsection{Proprietà formali}

SHAP identifica una classe di metodi denominati \textit{metodi di attribuzione additiva delle feature} e dimostra che esiste un'unica soluzione all'interno di questa classe che soddisfa simultaneamente tre proprietà fondamentali \cite{lundberg2017unified}:

\begin{itemize}
    \item \textbf{Accuratezza locale (\textit{Local Accuracy}).} La somma delle attribuzioni delle singole \textit{feature} ($\phi_i$) e del valore di base ($\phi_0$) deve essere uguale all'output del modello originale $f(x)$ per l'input considerato. Questa proprietà garantisce che la spiegazione sia completa: nessuna parte della previsione rimane non attribuita.
    \item \textbf{Assenza (\textit{Missingness}).} Se una \textit{feature} è assente dall'input originale, il suo impatto attribuito deve essere nullo. Questa proprietà assicura che non vengano assegnati contributi a informazioni non disponibili.
    \item \textbf{Consistenza (\textit{Consistency}).} Se un modello viene modificato in modo tale che il contributo marginale di una \textit{feature} aumenti o rimanga invariato, indipendentemente dagli altri input, l'attribuzione di quella \textit{feature} non deve diminuire. Questa proprietà impedisce che le spiegazioni si comportino in modo controintuitivo al variare del modello.
\end{itemize}

L'unicità della soluzione che soddisfa congiuntamente queste tre proprietà costituisce il principale argomento teorico a favore di SHAP rispetto ad altri metodi di attribuzione: le spiegazioni prodotte non sono una scelta arbitraria tra molte possibili, ma l'unica attribuzione coerente con i principi di equità ereditati dalla teoria dei giochi \cite{lundberg2017unified}.

\subsubsection{TreeExplainer: ottimizzazione per modelli ad albero}

Il calcolo esatto dei valori di Shapley richiede la valutazione di $2^M$ sottoinsiemi di \textit{feature}, dove $M$ è il numero totale di variabili, risultando computazionalmente proibitivo per la maggior parte delle applicazioni reali. Per i modelli basati su \textit{ensemble} di alberi decisionali --- fra cui XGBoost, impiegato nel caso studio della presente tesi --- è stato sviluppato l'algoritmo \textit{Tree SHAP}, che sfrutta la struttura degli alberi per calcolare i valori SHAP esatti in tempo polinomiale \cite{lundberg2019explainable}.

L'algoritmo utilizza due procedure ricorsive, denominate \textsc{Extend} e \textsc{Unwind}, che operano durante la discesa lungo l'albero. \textsc{Extend} tiene traccia, a ogni nodo, della proporzione di sottoinsiemi di \textit{feature} che fluiscono in ciascun ramo, distinguendo tra i casi in cui una \textit{feature} è presente (e segue il percorso dell'istanza osservata) e quelli in cui è assente (per i quali viene calcolata una media ponderata dei rami, pesata in base alla copertura dei dati di addestramento). \textsc{Unwind} interviene nei nodi foglia per annullare selettivamente l'effetto di ciascuna \textit{feature} e isolare il suo contributo specifico secondo i pesi di Shapley. Questo approccio consente di tracciare simultaneamente tutti i possibili sottoinsiemi in un unico passaggio attraverso l'albero, riducendo la complessità computazionale da $O(TLM \cdot 2^M)$ a $O(TLD^2)$, dove $T$ è il numero di alberi, $L$ il numero di foglie e $D$ la profondità massima \cite{lundberg2019explainable}.

Un aspetto rilevante riguarda il trattamento delle \textit{feature} assenti durante il calcolo. \textit{Tree SHAP} distingue due modalità di stima dell'aspettativa condizionata. La modalità \textit{tree path-dependent}, utilizzata per impostazione predefinita, sfrutta la copertura dei nodi dell'albero per pesare i rami quando una \textit{feature} è assente, rispecchiando la distribuzione dei dati appresa dal modello e tenendo conto implicitamente delle correlazioni tra le variabili. La modalità \textit{interventional} impone invece l'indipendenza tra le \textit{feature}, integrando l'output del modello su un insieme di campioni di riferimento. La prima modalità risulta più fedele al comportamento effettivo del modello sui dati reali, ma può attribuire importanza a \textit{feature} correlate anche in assenza di un impatto causale diretto; la seconda è più appropriata quando si desidera un'interpretazione di tipo causale \cite{lundberg2019explainable}.

\subsubsection{Dalle spiegazioni locali a quelle globali}

Sebbene SHAP sia concepito come un metodo di spiegazione locale, orientato alla singola previsione, l'aggregazione dei valori SHAP su un intero dataset consente di ricavare una visione globale dell'importanza delle \textit{feature} \cite{lundberg2019explainable}.

Lo strumento principale per questa transizione è il \textit{summary plot}, un grafico a sciame (\textit{beeswarm}) in cui ogni punto rappresenta un'istanza del dataset: la posizione sull'asse orizzontale indica l'impatto della \textit{feature} (il valore SHAP), mentre il colore codifica il valore della \textit{feature} stessa. Questa rappresentazione consente di visualizzare simultaneamente la magnitudo, la prevalenza e la direzione degli effetti di ciascuna variabile. Il \textit{dependence plot} offre un livello di dettaglio ulteriore, rappresentando il valore di una \textit{feature} sull'asse orizzontale e il corrispondente valore SHAP sull'asse verticale: la dispersione verticale tra punti con lo stesso valore della \textit{feature} rivela la presenza di interazioni con altre variabili \cite{lundberg2019explainable}.

\textit{Tree SHAP} consente inoltre di calcolare i \textit{valori di interazione SHAP}, basati sull'indice di interazione di Shapley. Questi valori producono, per ogni previsione, una matrice di attribuzioni in cui la diagonale contiene gli effetti principali di ciascuna \textit{feature} e gli elementi fuori diagonale rappresentano gli effetti di interazione puri tra coppie di variabili. Rispetto ai valori SHAP standard, che condensano effetto principale e interazioni in un unico numero, i valori di interazione consentono di separare le due componenti, rendendo visibili pattern che altrimenti rimarrebbero nascosti \cite{lundberg2019explainable}.

\subsubsection{Limiti}

Nonostante le solide basi teoriche, il framework SHAP presenta alcune criticità che è opportuno considerare. La principale riguarda la \textbf{complessità computazionale}: anche con le ottimizzazioni di \textit{Tree SHAP}, il calcolo dei valori di Shapley per modelli non arborei richiede approssimazioni basate su campionamento (\textit{Kernel SHAP}), che necessitano di un numero elevato di valutazioni per produrre stime stabili \cite{lundberg2017unified}. Un secondo limite concerne l'\textbf{assunzione di indipendenza delle \textit{feature}}: molte delle approssimazioni computazionali, tra cui \textit{Kernel SHAP}, presuppongono l'indipendenza o la linearità del modello per semplificare il calcolo delle aspettative condizionate. In presenza di \textit{feature} fortemente correlate, questa assunzione può condurre a spiegazioni che non riflettono fedelmente il comportamento del modello \cite{lundberg2017unified}. Infine, poiché la maggior parte dei modelli non è in grado di gestire pattern arbitrari di valori mancanti, SHAP deve \textbf{approssimare} l'assenza di una \textit{feature} integrando sulla sua distribuzione marginale, introducendo un livello di approssimazione rispetto al valore teorico puro \cite{lundberg2017unified}.

\subsection{LIME (Local Interpretable Model-agnostic Explanations)}

Il framework LIME (\textit{Local Interpretable Model-agnostic Explanations}), proposto da Ribeiro, Singh e Guestrin, adotta una strategia complementare rispetto a SHAP: anziché derivare le attribuzioni da un fondamento assiomatico, approssima localmente il comportamento di un modello complesso mediante un modello surrogato interpretabile, tipicamente lineare, addestrato nell'intorno dell'istanza da spiegare \cite{ribeiro2016why}.

\subsubsection{Perturbazione locale e modello surrogato lineare}

L'intuizione alla base di LIME è che, sebbene il confine decisionale di un modello complesso possa essere globalmente non lineare, in un intorno sufficientemente ristretto di una singola osservazione esso può essere approssimato con ragionevole fedeltà da un modello più semplice. Il metodo opera in tre fasi \cite{ribeiro2016why}.

Nella prima fase, LIME genera un insieme di campioni sintetici perturbati nell'intorno dell'istanza da spiegare. Per i dati tabulari, la perturbazione avviene estraendo valori da una distribuzione normale la cui media e deviazione standard sono ricavate dalla corrispondente \textit{feature} nel dataset di addestramento \cite{molnar2025interpretable}. Un aspetto rilevante di questa procedura è che il campionamento viene effettuato indipendentemente per ciascuna \textit{feature}, ignorando le eventuali correlazioni presenti nei dati reali: ciò può generare combinazioni di valori poco plausibili, un limite che il metodo condivide con altre tecniche di perturbazione \cite{molnar2025interpretable}.

Nella seconda fase, ciascun campione perturbato viene sottoposto al modello originale per ottenerne la previsione, che funge da etichetta per l'addestramento del surrogato. A ogni campione viene inoltre assegnato un peso proporzionale alla sua prossimità rispetto all'istanza di interesse, calcolato mediante un kernel esponenziale:
\[
\pi_x(z) = \exp\!\left(-\frac{D(x, z)^2}{\sigma^2}\right)
\]
dove $D$ è una funzione di distanza calcolata sui dati normalizzati e $\sigma$ è l'ampiezza del kernel, che determina la dimensione dell'intorno locale \cite{ribeiro2016why}. Campioni più vicini all'istanza originale ricevono un peso maggiore e influenzano in misura prevalente il modello surrogato, mentre campioni lontani contribuiscono in modo trascurabile.

Nella terza fase, LIME addestra un modello interpretabile $g$ --- tipicamente un modello lineare sparso --- sul dataset pesato di campioni perturbati, minimizzando la funzione obiettivo:
\[
\xi(x) = \underset{g \in G}{\text{argmin}} \; L(f, g, \pi_x) + \Omega(g)
\]
dove $L(f, g, \pi_x)$ misura l'infedeltà del surrogato $g$ rispetto al modello originale $f$ nella località definita da $\pi_x$, e $\Omega(g)$ è un termine di regolarizzazione che penalizza la complessità del modello, ad esempio il numero di \textit{feature} con coefficiente non nullo \cite{ribeiro2016why}. Il risultato è un insieme ristretto di \textit{feature} con i relativi pesi, che rappresentano l'importanza locale di ciascuna variabile per la previsione analizzata.

Un aspetto qualificante di LIME è la sua natura \textit{model-agnostic}: il metodo tratta il modello originale come una scatola nera, interagendo con esso esclusivamente attraverso coppie input-output, senza richiedere accesso alla struttura interna, ai gradienti o ai pesi del modello \cite{ribeiro2016why}. Questa proprietà garantisce l'applicabilità del framework a qualsiasi classe di modelli, ma comporta al contempo una dipendenza dalla qualità della rappresentazione interpretabile: se quest'ultima non è sufficientemente espressiva, LIME non sarà in grado di catturare aspetti rilevanti del comportamento del modello.

\subsubsection{Fedeltà locale}

La qualità di una spiegazione LIME è valutata in termini di \textit{fedeltà locale}, ovvero la capacità del modello surrogato di riprodurre fedelmente le previsioni del modello originale nell'intorno dell'istanza analizzata \cite{molnar2025interpretable}. Tale fedeltà è misurata dalla componente di perdita $L$ nella funzione obiettivo: un valore basso indica che il surrogato lineare approssima adeguatamente il comportamento locale della \textit{black box}, mentre un valore elevato segnala che l'approssimazione lineare è inadeguata in quel punto dello spazio delle \textit{feature} e che la spiegazione risultante è poco affidabile.

La fedeltà locale è intrinsecamente legata a due scelte progettuali che l'utente deve compiere prima di generare la spiegazione \cite{molnar2025interpretable}. La prima è l'\textbf{ampiezza del kernel} $\sigma$, che controlla la dimensione dell'intorno locale: un valore ridotto restringe l'analisi a un vicinato molto prossimo all'istanza, aumentando la fedeltà ma riducendo la quantità di informazione disponibile per il surrogato; un valore elevato amplia il vicinato, rischiando tuttavia di includere regioni dello spazio in cui il modello presenta un comportamento sostanzialmente diverso. La seconda è il \textbf{numero di \textit{feature}} $K$ ammesse nel surrogato: un valore di $K$ più alto consente al modello lineare di catturare dinamiche più complesse, accrescendo potenzialmente la fedeltà, ma al costo di una minore immediatezza interpretativa \cite{molnar2025interpretable}. Entrambe le scelte sono tipicamente definite in via euristica e la letteratura evidenzia come la determinazione dell'ampiezza ottimale del kernel costituisca un problema aperto \cite{molnar2025interpretable}.

\subsubsection{Limiti}

Il framework LIME presenta alcune criticità rilevanti che ne circoscrivono l'affidabilità e l'ambito di applicazione.

Il limite più discusso è l'\textbf{instabilità delle spiegazioni}. Poiché i campioni perturbati vengono generati tramite un processo stocastico, esecuzioni ripetute del metodo sulla medesima istanza possono produrre spiegazioni differenti, sia nel \textit{ranking} delle \textit{feature} sia nei valori dei coefficienti \cite{molnar2025interpretable}. Questa variabilità, particolarmente accentuata quando il numero di campioni è limitato o il confine decisionale locale presenta irregolarità, mina la riproducibilità delle spiegazioni e ne complica l'impiego in contesti decisionali dove la coerenza è un requisito.

Un secondo limite riguarda la \textbf{sensibilità alla parametrizzazione}. Come evidenziato nella sezione precedente, la scelta dell'ampiezza del kernel e del numero di \textit{feature} influenza in modo sostanziale il contenuto della spiegazione. In assenza di criteri formali per la calibrazione di questi parametri, valori diversi possono condurre a interpretazioni divergenti della medesima previsione \cite{molnar2025interpretable}.

Il terzo limite attiene alla natura stessa dell'\textbf{approssimazione lineare}. Se il modello originale presenta interazioni non lineari complesse anche in un intorno ristretto dell'istanza, un surrogato lineare non sarà in grado di catturarle fedelmente, producendo spiegazioni che semplificano eccessivamente il comportamento del modello \cite{molnar2025interpretable}.

Un ultimo limite, specifico per i dati tabulari, riguarda l'\textbf{assunzione di indipendenza} nel campionamento delle perturbazioni. Poiché i valori di ciascuna \textit{feature} vengono estratti indipendentemente dalla distribuzione marginale, i campioni sintetici possono rappresentare combinazioni di valori irrealistiche, con potenziali effetti sulla qualità delle previsioni ottenute dal modello originale e, di conseguenza, sulla fedeltà della spiegazione \cite{molnar2025interpretable}.

\subsection{Spiegazioni controfattuali (Alibi CounterfactualProto)}

Le spiegazioni controfattuali costituiscono un approccio concettualmente distinto rispetto a SHAP e LIME. Mentre questi ultimi tentano di ricostruire, con strumenti diversi, il contributo delle \textit{feature} alla previsione di un modello, il controfattuale rinuncia a spiegare la logica interna della \textit{black box} e si concentra su un'unica domanda: \textit{quali sono le modifiche minime da apportare all'istanza affinché la previsione cambi?} \cite{wachter2017counterfactual}. Il fondamento concettuale è la nozione di \textit{closest possible world}, ovvero il mondo possibile più vicino all'istanza osservata in cui il modello produrrebbe l'esito desiderato \cite{wachter2017counterfactual}. Tale formulazione è coerente con la tesi secondo cui le spiegazioni umane sono intrinsecamente \textit{contrastive}: non si spiega perché si è verificato un evento in assoluto, bensì perché esso si sia verificato al posto di un esito atteso alternativo \cite{miller2019explanation}.

Da questa impostazione discendono tre proprietà qualificanti. Il controfattuale è \textit{model-agnostic}, poiché interagisce con il modello originale esclusivamente attraverso coppie input-output, senza richiederne la struttura interna \cite{wachter2017counterfactual}. È \textit{locale}, poiché si riferisce alla singola decisione di un'istanza specifica e non al comportamento globale del modello. È \textit{orientato all'azione} (\textit{actionable}), poiché non si limita a descrivere ciò che il modello ha fatto, ma indica concretamente quali variabili il soggetto potrebbe modificare per conseguire un esito diverso in futuro \cite{wachter2017counterfactual}. Questa natura prescrittiva rende i controfattuali particolarmente rilevanti in contesti decisionali ad alto impatto, nei quali la possibilità per l'interessato di comprendere e contestare la decisione costituisce un requisito sostanziale \cite{wachter2017counterfactual}.

\subsubsection{Formulazione di Wachter et al.}

La formalizzazione più diffusa delle spiegazioni controfattuali è quella proposta da Wachter, Mittelstadt e Russell \cite{wachter2017counterfactual}. Dato un modello $f_w$, un'istanza originale $x_i$ e un target desiderato $y'$, il controfattuale $x'$ è la soluzione del problema di ottimizzazione:
\[
\arg\min_{x'} \max_{\lambda} \; \lambda \, \bigl(f_w(x') - y'\bigr)^2 + d(x_i, x')
\]
Il primo termine penalizza la distanza tra la previsione del controfattuale e il target desiderato, imponendo la \textit{validità} della soluzione; il secondo termine $d(x_i, x')$ è una funzione di distanza che misura quanto $x'$ si discosti dall'istanza originale, imponendo la \textit{prossimità} \cite{wachter2017counterfactual}. Il parametro $\lambda$ bilancia i due obiettivi ed è incrementato iterativamente fino al raggiungimento di una soluzione con predizione sufficientemente prossima a $y'$.

Una scelta progettuale rilevante riguarda la funzione di distanza $d$. Wachter et al. suggeriscono l'adozione della norma $L_1$, pesata per la deviazione assoluta mediana (MAD) delle singole \textit{feature}, in luogo della norma $L_2$ \cite{wachter2017counterfactual}. La motivazione è duplice. Da un lato, la ponderazione per MAD normalizza le variabili rispetto alla loro variabilità empirica, evitando che \textit{feature} con scale numeriche differenti dominino la misura. Dall'altro, la norma $L_1$ tende a produrre soluzioni \textit{sparse}, ovvero a concentrare i cambiamenti su un numero ridotto di variabili anziché distribuirli uniformemente su molte, come farebbe invece la norma $L_2$ \cite{wachter2017counterfactual}. La sparsity è considerata una proprietà qualificante per l'interpretabilità umana: la memoria di lavoro gestisce efficacemente solo un numero limitato di elementi simultaneamente, e una spiegazione che indichi poche leve d'intervento risulta nettamente più comprensibile di una che ne elenchi molte \cite{wachter2017counterfactual}.

\subsubsection{CounterfactualProto: guida tramite prototipi di classe}

Il metodo originario presenta un limite noto: ottimizzando soltanto validità e prossimità, esso tende a generare controfattuali che superano il confine decisionale del modello ma non necessariamente risiedono in regioni dello spazio delle \textit{feature} plausibili rispetto alla distribuzione dei dati reali \cite{vanlooveren2021counterfactualproto}. Il risultato può essere una spiegazione formalmente valida ma realisticamente implausibile, di scarsa utilità come guida all'azione.

Per affrontare questa criticità, Van Looveren e Klaise propongono \textit{CounterfactualProto}, un metodo \textit{model-agnostic} implementato nella libreria Alibi che introduce nel processo di ottimizzazione una guida esplicita verso la distribuzione della classe target \cite{vanlooveren2021counterfactualproto}. L'intuizione consiste nell'impiegare un \textbf{prototipo di classe}, ovvero una rappresentazione tipica delle istanze appartenenti alla classe desiderata, come attrattore che orienta le perturbazioni verso una regione \textit{in-distribution}.

La funzione di perdita di CounterfactualProto si compone di più termini, ciascuno associato a un obiettivo specifico \cite{vanlooveren2021counterfactualproto}. Un \textbf{termine di predizione} $L_{\text{pred}}$ incoraggia il modello ad associare all'istanza perturbata una classe diversa da quella originale. Una \textbf{distanza di tipo \textit{elastic net}}, combinazione delle norme $L_1$ e $L_2$, regola simultaneamente la sparsity e l'entità complessiva del cambiamento. Un \textbf{termine prototipo} $L_{\text{proto}}$ minimizza la distanza tra il controfattuale e il prototipo della classe target, forzando la soluzione verso una zona dello spazio delle \textit{feature} interpretabile e coerente con i dati osservati. Un ulteriore \textbf{termine \textit{autoencoder}} $L_{\text{AE}}$, opzionale, penalizza i controfattuali situati lontano dal \textit{manifold} dei dati di addestramento mediante l'errore di ricostruzione di un autoencoder.

La costruzione del prototipo può avvenire in due modi alternativi, a seconda della rappresentazione disponibile \cite{vanlooveren2021counterfactualproto}. Nella variante basata su \textbf{autoencoder}, il prototipo di ciascuna classe è definito come la codifica media, nello spazio latente dell'encoder, delle $K$ istanze di quella classe più vicine all'istanza originale; l'ottimizzazione avviene dunque in uno spazio compresso. Nella variante basata su \textbf{k-d tree}, in assenza di un encoder addestrato, per ciascuna classe viene costruito un k-d tree sulle istanze di addestramento e il prototipo è identificato come l'elemento più vicino all'istanza originale nello spazio delle \textit{feature} originarie. La prima variante si dimostra più efficace su dati ad alta dimensionalità come le immagini, mentre la seconda è adatta a dati tabulari o al caso in cui non si disponga di un autoencoder preaddestrato \cite{vanlooveren2021counterfactualproto}.

L'introduzione del termine prototipo apporta due vantaggi rispetto alla formulazione originaria. In primo luogo, migliora la \textit{plausibilità} delle spiegazioni: anziché limitarsi a superare il confine decisionale, l'ottimizzazione spinge attivamente la soluzione verso il centro della classe target, producendo controfattuali che somigliano a istanze tipiche e realistiche \cite{vanlooveren2021counterfactualproto}. In secondo luogo, accelera la \textit{convergenza} dell'ottimizzazione fornendo una direzione di ricerca informata; nei modelli \textit{black box}, in cui i gradienti del modello originale non sono direttamente disponibili e devono essere stimati numericamente, tale informazione aggiuntiva consente di omettere il termine $L_{\text{pred}}$ e riduce in misura rilevante il costo computazionale \cite{vanlooveren2021counterfactualproto}.

\subsubsection{Proprietà desiderabili}

La letteratura identifica un insieme di proprietà che un controfattuale di qualità dovrebbe soddisfare \cite{vanlooveren2021counterfactualproto, wachter2017counterfactual}. La \textbf{validità} richiede che la previsione del modello sull'istanza controfattuale corrisponda alla classe target. La \textbf{prossimità} richiede che il cambiamento rispetto all'istanza originale sia minimo, coerentemente con la nozione di \textit{mondo possibile più vicino}. La \textbf{sparsity} richiede che siano modificate poche \textit{feature}, per favorire la comprensibilità della spiegazione da parte di un utente non esperto. La \textbf{plausibilità}, o \textit{in-distribution-ness}, richiede che il controfattuale giaccia in prossimità della distribuzione dei dati di addestramento, e in particolare della sottodistribuzione della classe target. La \textbf{diversità} richiede che, quando esistono più percorsi controfattuali, ne venga restituito un insieme variegato, così da offrire all'utente alternative distinte. Infine, l'\textbf{actionability} richiede che i cambiamenti suggeriti riguardino variabili effettivamente modificabili dall'utente; in pratica ciò si ottiene vincolando durante l'ottimizzazione le variabili immutabili ai loro valori originari \cite{vanlooveren2021counterfactualproto}.

Queste proprietà non sono simultaneamente massimizzabili e la loro composizione richiede un compromesso, tipicamente governato dagli iperparametri della funzione di perdita.

\subsubsection{Limiti}

Le spiegazioni controfattuali presentano alcune criticità che ne circoscrivono l'affidabilità operativa.

Il primo limite, di natura strutturale, è la \textbf{non unicità}. Per una medesima istanza possono esistere numerosi controfattuali validi, che differiscono sia nelle \textit{feature} modificate sia nell'entità delle variazioni: la scelta dipende dalla funzione di distanza adottata e dagli iperparametri dell'ottimizzazione \cite{vanlooveren2021counterfactualproto, wachter2017counterfactual}. Questa molteplicità rende la nozione stessa di ``spiegazione corretta'' intrinsecamente soggettiva e può complicarne l'impiego in contesti nei quali è richiesta una giustificazione univoca.

Un secondo limite è la \textbf{sensibilità agli iperparametri}. CounterfactualProto dipende dalla calibrazione di diversi pesi --- tra cui il peso della sparsity, il peso del termine prototipo e il numero di vicini utilizzato per la costruzione del prototipo --- la cui scelta influenza in modo sostanziale l'esito della spiegazione; l'individuazione di valori appropriati è un'attività onerosa e priva di criteri formali generali \cite{vanlooveren2021counterfactualproto}.

Un terzo limite riguarda la \textbf{gestione di variabili categoriche e immutabili}. Per le \textit{feature} categoriche non esiste una nozione naturale di distanza, e CounterfactualProto ricorre a rappresentazioni \textit{embedding} basate su metriche specifiche, al costo di un'ulteriore complessità metodologica \cite{vanlooveren2021counterfactualproto}. Gli algoritmi standard, inoltre, non distinguono autonomamente tra variabili modificabili e variabili che non lo sono per ragioni logiche, etiche o legali: in assenza di vincoli espliciti il metodo può generare controfattuali che suggeriscono modifiche a variabili quali età, genere o etnia, producendo spiegazioni inutilizzabili e potenzialmente discriminatorie \cite{wachter2017counterfactual, vanlooveren2021counterfactualproto}.

Un quarto limite concerne il rischio di \textbf{controfattuali non azionabili} per ragioni causali. La formulazione standard assume implicitamente l'indipendenza delle variabili, mentre nella realtà il cambiamento di una \textit{feature} può presupporre o comportare il cambiamento di altre a essa causalmente collegate: in assenza di un modello causale esplicito, il controfattuale può suggerire modifiche che, pur minimizzando la funzione di perdita, risultano logicamente o praticamente irrealizzabili \cite{wachter2017counterfactual}.

Un ultimo limite è specifico di CounterfactualProto e consiste in un \textbf{trade-off tra sparsity e plausibilità}. La guida verso il prototipo di classe avvicina la soluzione a istanze tipiche della classe target e ne aumenta la plausibilità, ma tende a produrre controfattuali meno sparsi rispetto a quelli ottenuti minimizzando unicamente la distanza $L_1$: l'obiettivo di somigliare a un esempio reale della classe target richiede, in genere, l'aggiustamento di più \textit{feature} simultaneamente \cite{vanlooveren2021counterfactualproto}. La scelta tra massimizzare la sparsity o la plausibilità è dunque una decisione progettuale che dipende dal contesto applicativo.

\subsection{Confronto teorico SHAP e LIME}

SHAP e LIME condividono l'obiettivo di attribuire a ciascuna \textit{feature} un peso che rifletta il suo contributo a una specifica previsione, ma vi giungono per vie metodologicamente distinte: il primo deriva le proprie attribuzioni da un quadro assiomatico radicato nella teoria dei giochi cooperativi, mentre il secondo costruisce empiricamente un surrogato lineare locale nell'intorno dell'istanza da spiegare \cite{lundberg2017unified, ribeiro2016why}. Le conseguenze di questa differenza si riflettono su più dimensioni operative, che la Tabella~\ref{tab:shap-lime} sintetizza in forma sinottica lungo i sei assi individuati come rilevanti nelle sezioni precedenti.

\begin{table}[htbp]
\centering
\small
\caption{Confronto sinottico tra SHAP e LIME lungo sei dimensioni rilevanti per l'applicazione pratica.}
\label{tab:shap-lime}
\begin{tabular}{@{}p{2.6cm}p{5.3cm}p{5.3cm}@{}}
\toprule
\textbf{Dimensione} & \textbf{SHAP} & \textbf{LIME} \\
\midrule
Ambito & Nativamente locale, con un'aggregazione verso una vista globale supportata dal framework attraverso \textit{summary plot} e valori di interazione. & Strettamente locale; una vista globale è ottenibile soltanto tramite aggregazione ex post (ad esempio la media del valore assoluto dei pesi) priva di un fondamento teorico analogo. \\
\addlinespace
Fedeltà & Fedeltà locale esatta per costruzione: l'assioma di \textit{Local Accuracy} impone $\sum_i \phi_i + \phi_0 = f(x)$. & Fedeltà locale approssimata, misurata dalla componente di perdita $L$; degrada quando il confine decisionale è non lineare anche nell'intorno. \\
\addlinespace
Stabilità & \textit{Tree SHAP} è deterministico, poiché calcola i valori di Shapley in forma esatta; \textit{Kernel SHAP}, fondato sul campionamento, condivide la stessa instabilità di LIME. & Stocastico per natura: esecuzioni ripetute sulla medesima istanza possono produrre \textit{ranking} e coefficienti differenti. \\
\addlinespace
Velocità & \textit{Tree SHAP} è polinomiale ($O(TLD^2)$) e mantiene tempi pressoché costanti al crescere del numero di \textit{feature}, risultando significativamente più efficiente di LIME sui modelli ad albero. & Costo lineare nel numero di \textit{feature} e dominato da $N$ valutazioni del modello originale; la variabilità tra classi di modelli è trascurabile, ma il numero di inferenze richieste rende il metodo oneroso in alta dimensionalità. \\
\addlinespace
Interpretabilità & Attribuzione additiva, con valori $\phi_i$ espressi nella medesima unità dell'output del modello; la scomposizione della previsione è esaustiva. & Coefficienti di un modello lineare sparso definito sulla rappresentazione interpretabile; leggibili come regole semplici, ma privi di garanzie formali di completezza. \\
\addlinespace
Dipendenza dal modello & Offre varianti \textit{model-agnostic} (\textit{Kernel SHAP}) e varianti \textit{model-specific} ottimizzate (\textit{Tree SHAP} per \textit{ensemble} di alberi). & Pienamente \textit{model-agnostic}: il modello è interrogato esclusivamente come oracolo input-output, senza alcuna assunzione sulla sua struttura interna. \\
\bottomrule
\end{tabular}
\end{table}

Le evidenze empiriche raccolte in un'analisi comparativa dei metodi additivi di spiegazione locale confermano il quadro delineato sul versante computazionale: per i modelli ad albero, \textit{Tree SHAP} risulta nettamente più efficiente di LIME sia in bassa sia in alta dimensionalità, mentre il costo di LIME si mantiene pressoché indipendente dalla classe del modello sottostante ma cresce linearmente con il numero di \textit{feature} e con il numero di campioni perturbati richiesti per ottenere spiegazioni stabili \cite{doumard2022comparative}.

Sul piano pratico, le due famiglie rispondono dunque a esigenze parzialmente diverse. Nel contesto della presente tesi, in cui il modello predittivo è un \textit{ensemble} di alberi (XGBoost), \textit{Tree SHAP} offre attribuzioni esatte, deterministiche e computazionalmente trattabili, e costituisce il riferimento metodologico primario per le spiegazioni locali. LIME svolge in questo quadro una funzione complementare: la sua natura \textit{model-agnostic} e la forma lineare del surrogato forniscono una seconda lettura della medesima previsione, utile a valutarne la robustezza e a mettere in luce eventuali discordanze, riconducibili alla pluralità di letture compatibili con uno stesso modello già evocata discutendo l'effetto Rashomon.
